% SGA 3, Expose I, Loop 1 English reconstruction.
% Scope: Section 6, "G-equivariant objects and modules," through the
% bibliography and end of Expose I.
% Authority: Polo--Gille Exp1-13oct24.pdf, component-local pages 39--48
% (combined-reader pages 49--58). Force-OCR and the English Markdown by
% Jacob Reinhold are drafting and comparison witnesses only.

% Source page: Exp1 component-local p. 39; combined-reader p. 49.
\section*{\SGAInlineRefTarget{sga3:i:section:6}6. \(G\)-等变对象和模}
\addcontentsline{toc}{section}{6. G-等变对象和模}

设 \(\mathscr C\) 是一个具有终对象且存在纤维积的范畴.
\footnote{编者注: 本节为新增内容.}
设 \(G\) 是一个 \(\widehat{\mathscr C}\)-群,
\(\pi:M\to X\) 是 \(\widehat{\mathscr C}\) 中的一个态射,
\(\lambda=\lambda_X:G\times X\to X\) 是 \(G\) 在 \(X\) 上的一个作用.
以下以 \(Y\times_f M\) 表示 \(\pi:M\to X\) 与 \(X\)-函子
\(f:Y\to X\) 的纤维积.

对每个 \(U\in\operatorname{Ob}(\mathscr C)\) 和 \(x\in X(U)\), 记
\(M_x=U\times_x M\); 因而对每个 \(\phi:U'\to U\), 有
\[
  M_x(U')
  =\{m\in M(U')\mid \pi(m)=x_{U'}=\phi^*(x)\}.
\]
最后, 若 \(g\in G(U)\), 我们也以 \(gx\) 表示 \(X(U)\) 的元素
\(\lambda(g,x)\).

\medskip
\noindent\SGAInlineRefTarget{sga3:i:definition:6-1}\textbf{定义 6.1.}
\textit{\SGAInlineRefTarget{sga3:i:definition:6-1:item:a}(a) 若给定了 \(G\) 在 \(M\) 上提升 \(\lambda\) 的作用
\(\Lambda:G\times M\to M\), 就称 \(M\) 是一个 \(G\)-等变（equivariant）\(X\)-对象;
也就是说, 下列方块交换.}

\begin{figure}[H]
  \centering
  \begin{tikzcd}[column sep=4.7em,row sep=3.2em]
  G\times M
    \arrow[r,"\Lambda"]
    \arrow[d,"\operatorname{id}_G\times\pi"']
  &
  M
    \arrow[d,"\pi"]
  \\
  G\times X
    \arrow[r,"\lambda"]
  &
  X
\end{tikzcd}

  \phantomsection\label{fig:sga3-exp1-p039-equivariant-action-square}
\end{figure}

\textit{等价地, 对每个态射 \((g,x):U\to G\times X\), 给定映射}
\[
  \Lambda_x^U(g):M_x(U)\longrightarrow M_{gx}(U),
  \qquad m\longmapsto g\cdot m,
\]
\textit{满足 \(1\cdot m=m\) 和
\(g\cdot(h\cdot m)=(gh)\cdot m\), 且关于 \((G\times X)\)-对象
\(U\) 具有函子性. 另一种等价的说法是, 给定 \(U\)-对象的态射}
\[
  \Lambda_x(g):M_x\longrightarrow M_{gx}
\]
\textit{满足
\(\Lambda_x(1)=\operatorname{id}\) 和
\(\Lambda_{hx}(g)\circ\Lambda_x(h)=\Lambda_x(gh)\).}

\textit{(b) 现设 \(\mathbf O\) 是一个
\(\widehat{\mathscr C}\)-环, 并置
\(\mathbf O_X=\mathbf O\times X\). 在~\SGARefLink{sga3:i:definition:6-1:item:a}{(a)} 的假设下,
若 \(M\) 是一个 \(\mathbf O_X\)-模（参见\SGARefLink{sga3:i:paragraph:4-3-3-1}{定义~4.3.3.1},
该定义对范畴 \(\mathscr C\) 上的任意环函子均适用）, 且作用 \(\Lambda\)
与 \(M\) 上的 \(\mathbf O_X\)-模结构相容, 就称之为一个
\(G\)-等变 \(\mathbf O_X\)-模. 因而, 对每个
\((g,x)\in G(U)\times X(U)\),
\(U\in\operatorname{Ob}(\mathscr C)\), 映射}
\[
  \Lambda_x(g):M_x\longrightarrow M_{gx},
  \qquad m\longmapsto g\cdot m,
\]
\textit{都是 \(\mathbf O_U\)-模的态射.}

\medskip
\noindent\SGAInlineRefTarget{sga3:i:remark:6-2}\textbf{注 6.2.}
(i) 在~\SGARefLink{sga3:i:definition:6-1:item:a}{(a)} 中, 条件
\(\Lambda_x(1)=\operatorname{id}\) 和
\(\Lambda_{hx}(g)\circ\Lambda_x(h)=\Lambda_x(gh)\) 显然蕴含每个
\(\Lambda_x(g)\) 都是同构, 其逆为
\(\Lambda_{gx}(g^{-1})\). 反之, 若每个 \(\Lambda_x(g)\) 都是同构,
在第二个条件中取 \(h=1\), 就得到
\(\Lambda_x(1)=\operatorname{id}\).

(ii) 设 \(M\) 是一个 \(\mathbf O_X\)-模. 给定一个使\SGARefLink{sga3:i:definition:6-1}{定义~6.1}
中的图交换的态射 \(\Lambda:G\times M\to M\), 并要求每个态射

% Source page: Exp1 component-local p. 40; combined-reader p. 50.
\[
  \Lambda_x(g):M_x\longrightarrow M_{gx},
  \qquad m\longmapsto g\cdot m,
\]
都是 \(\mathbf O_U\)-模的同构, 等价于给定一个
\(\mathbf O_{G\times X}\)-模的同构
\[
  \theta:
  G\times M=(G\times X)\times_{\mathrm{pr}_X}M
  \xrightarrow{\ \sim\ }(G\times X)\times_\lambda M,
  \qquad
  (g,x,m)\longmapsto(g,x,g\cdot m).
\]
由于已假设每个 \(\Lambda_x(g)\) 都是同构, 在
\(\Lambda_{hx}(g)\circ\Lambda_x(h)=\Lambda_x(gh)\) 中取 \(h=1\),
就得到等式 \(\Lambda_x(1)=\operatorname{id}\).
因此, \(\Lambda\) 是 \(G\) 在 \(M\) 上的作用, 即满足
\(g\cdot(h\cdot m)=(gh)\cdot m\), 当且仅当下列
\((G\times G\times X)\)-同构图交换. 这里 \(\mu\) 表示 \(G\) 的乘法,
\(f^*(\theta)\) 是由 \(\theta\) 沿
\(f:G\times G\times X\to G\times X\) 作基变换而得到的.

\begin{figure}[H]
  \centering
  \begin{tikzcd}[column sep=1.7em,row sep=3.2em]
  (G\times G\times X)
    \mathop{\times}\limits_{\mathrm{pr}_X\circ\mathrm{pr}_{2,3}} M
    \arrow[rr,"\mathrm{pr}_{2,3}^{*}(\theta)","\sim"']
    \arrow[d,equal]
  &&
  (G\times G\times X)
    \mathop{\times}\limits_{\lambda\circ\mathrm{pr}_{2,3}} M
    \arrow[d,equal]
  \\
  (G\times G\times X)
    \mathop{\times}\limits_{\mathrm{pr}_X\circ(\mu\times\operatorname{id}_X)} M
    \arrow[d,"(\mu\times\operatorname{id}_X)^*(\theta)"',"\sim"]
  &&
  (G\times G\times X)
    \mathop{\times}\limits_{\mathrm{pr}_X\circ(\operatorname{id}_G\times\lambda)} M
    \arrow[d,"(\operatorname{id}_G\times\lambda)^*(\theta)","\sim"']
  \\
  (G\times G\times X)
    \mathop{\times}\limits_{\lambda\circ(\mu\times\operatorname{id}_X)} M
    \arrow[rr,equal]
  &&
  (G\times G\times X)
    \mathop{\times}\limits_{\lambda\circ(\operatorname{id}_G\times\lambda)} M
  \qquad .
\end{tikzcd}

  \phantomsection\label{fig:sga3-exp1-p040-equivariance-cocycle}
\end{figure}

所以, 在 \(M\) 上给定 \(G\)-等变 \(\mathbf O_X\)-模结构,
等价于给定一个如上所述的 \(\mathbf O_{G\times X}\)-模同构
\(\theta\), 并使此图交换.

(iii) 以上各点都可推广到 \(G\) 仅为一个
\(\widehat{\mathscr C}\)-幺半群的情形. 此时, 给定提升 \(\lambda\)
的作用 \(\Lambda:G\times M\to M\), 并要求每个
\(\Lambda_x(g):M_x\to M_{gx}\) 都是 \(\mathbf O_U\)-模的态射,
等价于给定一个如~(ii) 中的 \(\mathbf O_{G\times X}\)-模态射
\(\theta\), 使得将上图各箭头下的同构符号去掉后该图交换, 且满足
\[
  p_M\circ\theta\circ(\varepsilon_G\times\operatorname{id}_M)
  =\operatorname{id}_M,
\]
其中 \(\varepsilon_G\) 是 \(G\) 的单位截面, \(p_M\) 是到 \(M\) 的投影.

\subsection*{\SGAInlineRefTarget{sga3:i:subsection:6-3}6.3. \(G\)-等变态射}

设 \(Y\) 是 \(\widehat{\mathscr C}\) 的另一个对象, 其上赋予了作用
\(\lambda_Y:G\times Y\to Y\); 又设 \(N\) 是另一个
\(G\)-等变 \(\mathbf O_X\)-模. \(\widehat{\mathscr C}\) 中的态射
\(f:Y\to X\)（分别地, \(\mathbf O_X\)-模的态射 \(\phi:M\to N\)）
若与 \(G\) 的作用交换, 就称为 \emph{\(G\)-等变的}.
用集合论记号表示, 这就是说
\[
  f(g\cdot y)=g\cdot f(y)
  \qquad
  \bigl(\text{分别地, }\phi(g\cdot m)=g\cdot\phi(m)\bigr),
\]
或者等价地,
\[
  f\circ\lambda_Y=\lambda_X\circ(\operatorname{id}_G\times f),
  \qquad
  \phi\circ\Lambda_M=\Lambda_N\circ(\operatorname{id}_G\times\phi).
\]
下列引理是直接的.

% Source page: Exp1 component-local p. 41; combined-reader p. 51.
\medskip
\noindent\SGAInlineRefTarget{sga3:i:lemma:6-3-1}\textbf{引理 6.3.1.}
\textit{设 \(f:Y\to X\) 是一个 \(G\)-等变态射, \(M\) 是一个
\(G\)-等变 \(\mathbf O_X\)-模. 则逆像
\(f^*(M)=Y\times_fM\) 是一个 \(G\)-等变
\(\mathbf O_Y\)-模.}

另一方面, \(G\) 作用在
\(\operatorname{Hom}_{\widehat{\mathscr C}}(Y,X)\) 上. 事实上, 设
\(T\in\operatorname{Ob}(\mathscr C)\), \(g\in G(T)\), 且 \(\phi\)
是一个 \(T\)-态射 \(Y_T\to X_T\). 则 \(g\) 给出 \(Y_T\) 和
\(X_T\) 的自同构 \(\lambda_Y(g)\) 和 \(\lambda_X(g)\), 我们以
\(g\cdot\phi\) 或 \(g\phi g^{-1}\) 表示态射
\[
  \lambda_X(g)\circ\phi\circ\lambda_Y(g^{-1}).
\]
这就给出了 \(G(T)\) 在
\(\operatorname{Hom}_T(Y_T,X_T)\) 上的作用, 且关于 \(T\) 具有函子性.

\medskip
\noindent\SGAInlineRefTarget{sga3:i:definition:6-3-2}\textbf{定义 6.3.2.}
\textit{因此, 对 \(\widehat{\mathscr C}\) 中任意态射 \(\phi:Y\to X\),
可以考虑它的稳定子 \(\operatorname{Stab}_G(\phi)\)
（参见 \SGARefLink{sga3:i:scholium:2-3-3-1}{2.3.3.1}）. 对每个
\(T\in\operatorname{Ob}(\mathscr C)\),
\(\operatorname{Stab}_G(\phi)(T)\) 是 \(G(T)\) 中所有满足
\(g\circ\phi_T=\phi_T\circ g\) 的 \(g\) 所组成的子群,
也就是说, 它们使下图交换.}

\begin{figure}[H]
  \centering
  \begin{tikzcd}[column sep=4.0em,row sep=3.0em]
  Y_T
    \arrow[r,"\phi_T"]
    \arrow[d,"g"']
  &
  X_T
    \arrow[d,"g"]
  \\
  Y_T
    \arrow[r,"\phi_T"]
  &
  X_T
\end{tikzcd}

  \phantomsection\label{fig:sga3-exp1-p041-stabilizer-square}
\end{figure}

\textit{于是 \(\phi:Y\to X\) 关于 \(\operatorname{Stab}_G(\phi)\) 等变.}

\subsection*{\SGAInlineRefTarget{sga3:i:subsection:6-4}6.4. 整体截面}

设 \(M\) 是一个 \(G\)-等变 \(\mathbf O_X\)-模.
以 \(S_0\) 表示 \(\mathscr C\) 的终对象, 并以
\(\prod_{X/S_0}M\) 表示\emph{\(M\) 在 \(X\) 上的截面函子}
（参见\SGARefLink{sga3:ii:definition:1-1}{第~II 篇, 1.1}）:
它将每个 \(T\in\operatorname{Ob}(\mathscr C)\) 映为
\[
  \operatorname{Hom}_X(X_T,M)
  =\operatorname{Hom}_{X_T}(X_T,M_T)
  =\Gamma(M_T/X_T).
\]
还要回顾, \(X\) 上的 \(\widehat{\mathscr C}\)-对象的每个态射
\(g:Z\to Y\) 都诱导一个阿贝尔群态射
\[
  M(g):M(Y)\longrightarrow M(Z),
\]
它与环同态 \(g^*:\mathbf O(Y)\to\mathbf
O(Z)\) 相容. 特别地, 当 \(Z=Y\), 即 \(g\) 是 \(Y\) 的一个
\(X\)-自同态时, 就得到阿贝尔群态射
\[
  M(g):M(Y)\longrightarrow M(Y)
\]
它一般不是 \(\mathbf O(Y)\)-模的态射, 但对每个
\(m\in M(Y)\) 和 \(\alpha\in\mathbf O(Y)\), 满足
\[
  M(g)(\alpha\cdot m)=g^*(\alpha)\cdot M(g)(m)
\]
. 以下简单地以 \(g^*\) 代替 \(M(g)\).

设 \(T\in\operatorname{Ob}(\mathscr C)\), 置 \(X_T=X\times T\),
并设 \(\mathrm{pr}_X:X_T\to X\) 是投影.
对 \(\alpha\in\mathbf O(X_T)\) 和 \(g\in G(T)\), 置
\[
  g(\alpha)=(g^{-1})^*(\alpha).
\]
用集合论语言来说, \(\mathbf O(X_T)\) 的这个元素由如下等式定义:
对每个 \(S\in\operatorname{Ob}(\mathscr C)\), \(x\in X(S)\) 和
\(t\in T(S)\), 有
\[
  g(\alpha)(x,t)=\alpha(g^{-1}x,t).
\]
这给出了 \(G(T)\) 通过环自同构在 \(\mathbf O(X_T)\) 上的左作用,
且关于 \(T\) 具有函子性; 当 \(\alpha\) 是 \(\mathbf O(T)\)
中某个元素在 \(\mathbf O(X_T)\) 中的像时, 有 \(g(\alpha)=\alpha\).

现设 \(\phi\) 为 \(X_T\) 的恒等态射（参见 \SGARefLink{sga3:i:subsection:6-3}{6.3}
以及下面对态射 \(\phi:Y\to X\) 的推广）.
根据分别是通过 \(\mathrm{pr}_X\circ\lambda(g^{-1})_T\) 还是通过
\(\mathrm{pr}_X\) 将 \(X_T\) 视为 \(X\)-对象, 将
\(\operatorname{Hom}_X(X_T,M)\) 分别记为 \(M(\phi g^{-1})\) 或 \(M(\phi)\).

% Source page: Exp1 component-local p. 42; combined-reader p. 52.
于是 \(\lambda(g^{-1})_T\) 是这两个 \(X\)-对象之间的 \(X\)-态射.
因而有阿贝尔群态射
\[
  (g^{-1})^*:M(\phi)\longrightarrow M(\phi g^{-1}),
  \qquad
  m\longmapsto m\circ\lambda(g^{-1})_T,
\]
且对 \(m\in M(\phi)\) 和 \(\alpha\in\mathbf O(X_T)\), 有
\[
  (g^{-1})^*(\alpha\cdot m)
  =(g\alpha)\cdot(g^{-1})^*(m)
\]
. 若 \(m\) 是 \(M_T\) 在 \(X_T\) 上的截面, 则
\((g^{-1})^*(m)\) 是用集合论方式由
\((x,t)\mapsto m(g^{-1}x,t)\) 定义的截面.
特别地, \((g^{-1})^*\) 是 \(\mathbf O(T)\)-模的态射.

\(\mathbf O(X_T)\)-模态射 \(\Lambda_x(g)\) 的函子性给出下列交换图.

\begin{figure}[H]
  \centering
  \begin{tikzcd}[column sep=4.5em,row sep=3.3em]
  M(\phi)
    \arrow[r,"(g^{-1})^*"]
    \arrow[d,"\Lambda_\phi(g)"']
  &
  M(\phi g^{-1})
    \arrow[d,"\Lambda_{\phi g^{-1}}(g)"]
  \\
  M(g\phi)
    \arrow[r,"(g^{-1})^*"]
  &
  M(g\phi g^{-1})
\end{tikzcd}

  \phantomsection\label{fig:sga3-exp1-p042-section-naturality-square}
\end{figure}

此外, 由于 \(\phi\) 是恒等态射, 有 \(g\phi g^{-1}=\phi\).
记
\[
  A(g)=(g^{-1})^*\circ\Lambda_\phi(g),
\]
则得到阿贝尔群态射
\[
  A(g):M(\phi)\longrightarrow M(\phi)
\]
它与 \(G(T)\) 在 \(\mathbf O(X_T)\) 上的作用相容, 即
\[
  A(g)(\alpha\cdot m)=(g\alpha)\cdot A(g)(m).
\]

最后, 若 \(h\) 是 \(G(T)\) 的另一个元素, 则 \(M\) 及态射
\(\Lambda_x(g)\) 的函子性给出下列交换图.

\begin{figure}[H]
  \centering
  \begin{tikzcd}[column sep=4.5em,row sep=3.3em]
  M(\phi)
    \arrow[d,"\Lambda_\phi(g)"']
    \arrow[dr,"\Lambda_\phi(hg)"]
  &
  \\
  M(g\phi)
    \arrow[r,"\Lambda_\phi(h)"]
    \arrow[d,"(g^{-1})^*"']
  &
  M(hg\phi)
    \arrow[d,"(g^{-1})^*"']
    \arrow[dr,"((hg)^{-1})^*"]
  &
  \\
  M(\phi)
    \arrow[r,"\Lambda_\phi(h)"]
  &
  M(h\phi)
    \arrow[r,"(h^{-1})^*"]
  &
  M(\phi)
\end{tikzcd}

  \phantomsection\label{fig:sga3-exp1-p042-section-action-composition}
\end{figure}

由此得 \(A(h)\circ A(g)=A(hg)\), 因而有:

\medskip
\noindent\SGAInlineRefTarget{sga3:i:proposition:6-4-1}\textbf{命题 6.4.1.}
\textit{对每个 \(T\), \(\mathbf O(X_T)\)-模
\[
  \left(\prod_{X/S_0}M\right)(T)=\Gamma(M_T/X_T)
\]
都赋有关于 \(T\) 具函子性的 \(G(T)\) 的作用, 它与 \(G(T)\) 在
\(\mathbf O(X_T)\) 上的作用相容; 也就是说,}
\[
  A(g)(\alpha\cdot m)=g(\alpha)\cdot A(g)(m).
\]

由于对 \(\alpha\in\mathbf O(T)\) 有 \(g(\alpha)=\alpha\),
特别地, \(\prod_{X/S_0}M\) 是一个
\(G\)-\(\mathbf O_{S_0}\)-模.

更一般地, 设 \(Y\) 是另一个 \(G\)-等变
\(\widehat{\mathscr C}\)-对象, \(\phi:Y\to X\) 是
\(\widehat{\mathscr C}\) 中的一个态射, 并如~\SGARefLink{sga3:i:subsection:6-3}{6.3}
中那样置 \(H=\operatorname{Stab}_G(\phi)\). 则
\[
  M_\phi=Y\times_\phi M
\]
是一个 \(H\)-等变 \(\mathbf O_Y\)-模. 我们得到:

% Source page: Exp1 component-local p. 43; combined-reader p. 53.
\medskip
\noindent\SGAInlineRefTarget{sga3:i:corollary:6-4-2}\textbf{推论 6.4.2.}
\textit{函子 \(\prod_{Y/S_0}M_\phi\) 是一个
\(\operatorname{Stab}_G(\phi)\)-\(\mathbf O_{S_0}\)-模.}

\subsection*{\SGAInlineRefTarget{sga3:i:subsection:6-5}6.5. \(G\)-等变 \(\mathcal O_X\)-模}

将前面的讨论用于下述情形. 设 \(S\) 是一个概形,
\(G\) 是一个作用在 \(S\)-概形 \(X\) 上的 \(S\)-群概形,
\(\mathcal F\) 是一个未必拟凝聚的 \(\mathcal O_X\)-模.

\medskip
\noindent\SGAInlineRefTarget{sga3:i:definition:6-5-1}\textbf{定义 6.5.1.}
\textit{若 \(\mathbf O_X\)-模 \(M=W(\mathcal F)\) 是
\(G\)-等变的, 就称 \(\mathcal F\) 是一个 \(G\)-等变
\(\mathcal O_X\)-模. 因而, 对每个态射
\((g,x):U\to G\times_S X\), 给定 \(\mathcal O_U\)-模的同构}
\[
  \Lambda_x(g):
  W\bigl(x^*(\mathcal F)\bigr)
  \xrightarrow{\ \sim\ }
  W\bigl((gx)^*(\mathcal F)\bigr),
\]
\textit{它们关于 \(U\) 具有函子性, 且满足
\(\Lambda_{hx}(g)\circ\Lambda_x(h)=\Lambda_x(gh)\).}
由于 \(W\) 是全忠实的（参见 \SGARefLink{sga3:i:remark:4-6-1-1}{4.6.1.1}
和 \SGARefLink{sga3:i:proposition:4-6-2}{4.6.2}）, 由此得到
\(\mathcal O_U\)-模的同构
\[
  \Lambda_x(g):x^*(\mathcal F)
  \xrightarrow{\ \sim\ }(gx)^*(\mathcal F),
\]
其中 \(gx\) 表示 \(\lambda\circ(g,x):U\to X\).
将这一点用于 \(G\times_S X\) 的恒等态射, 就得到
\(\mathcal O_{G\times_S X}\)-模的同构
\[
  (\star)\qquad
  \theta:\mathrm{pr}_X^*(\mathcal F)
  \xrightarrow{\ \sim\ }\lambda^*(\mathcal F)
\]
它使下列 \(G\times_SG\times_SX\) 上的层态射图 \((\star\star)\) 交换.

\begin{figure}[H]
  \centering
  \begin{tikzcd}[column sep=1.8em,row sep=3.5em]
  \mathrm{pr}_{2,3}^{*}\circ\mathrm{pr}_X^{*}(\mathcal F)
    \arrow[rr,"\mathrm{pr}_{2,3}^{*}(\theta)","\sim"']
    \arrow[d,equal]
  &&
  \mathrm{pr}_{2,3}^{*}\circ\lambda^{*}(\mathcal F)
    \arrow[rr,equal]
  &&
  (\operatorname{id}_G\times\lambda)^{*}\circ\mathrm{pr}_X^{*}(\mathcal F)
    \arrow[d,"(\operatorname{id}_G\times\lambda)^*(\theta)","\sim"']
  \\
  (\mu\times\operatorname{id}_X)^{*}\circ\mathrm{pr}_X^{*}(\mathcal F)
    \arrow[rr,"(\mu\times\operatorname{id}_X)^*(\theta)","\sim"']
  &&
  (\mu\times\operatorname{id}_X)^{*}\circ\lambda^{*}(\mathcal F)
    \arrow[rr,equal]
  &&
  (\operatorname{id}_G\times\lambda)^{*}\circ\lambda^{*}(\mathcal F)
  \qquad .
\end{tikzcd}

  \phantomsection\label{fig:sga3-exp1-p043-sheaf-cocycle}
\end{figure}

若 \(\mathcal E\) 是另一个 \(G\)-等变 \(\mathcal O_X\)-模,
则 \(\mathcal O_X\)-模的态射 \(\phi:\mathcal E\to\mathcal F\)
若使 \(W(\phi):W(\mathcal E)\to W(\mathcal F)\) 为 \(G\)-等变的,
就称为\emph{\(G\)-等变的}. 用上述记号表示, 这等价于
\[
  \theta_{\mathcal F}\circ\mathrm{pr}_X^*(\phi)
  =\lambda^*(\phi)\circ\theta_{\mathcal E}.
\]

\medskip
\noindent\SGAInlineRefTarget{sga3:i:remark:6-5-2}\textbf{注 6.5.2.}
若 \(f:Y\to X\) 是 \(G\)-等变的, 则\SGARefLink{sga3:i:lemma:6-3-1}{引理~6.3.1}
表明 \(f^*(\mathcal F)\) 是一个 \(G\)-等变
\(\mathcal O_Y\)-模.

\medskip
\noindent\SGAInlineRefTarget{sga3:i:remark:6-5-3}\textbf{注 6.5.3.}
再假设 \(\mathcal F\) 是拟凝聚的. 此时, 在
\(M=W(\mathcal F)\) 上给定 \(G\)-等变 \(\mathcal O_X\)-模结构,
等价于在 \(V(\mathcal F)\) 上给定这样的结构, 而这又等价于给定
\(G\) 在向量纤维化 \(V(\mathcal F)\) 上的作用, 它与 \(X\) 上的作用相容,
并在各纤维上是线性的.

事实上, 设
\[
  \varphi:\lambda^*(\mathcal F)
  \xrightarrow{\ \sim\ }\mathrm{pr}_X^*(\mathcal F)
\]
是 \(\theta\) 的逆. 对每个态射
\((g,x):U\to G\times_S X\), 从 \((gx)^*(\mathcal F)\) 到
\(x^*(\mathcal F)\) 的 \(\mathcal O_U\)-模同构
\[
  \varphi_x(g)
  =\Lambda_x(g)^{-1}
  =\Lambda_{gx}(g^{-1})
\]
诱导一个 \(\mathcal O_U\)-模同构
\[
  V\bigl((gx)^*(\mathcal F)\bigr)
  \xrightarrow{\ \sim\ }V\bigl(x^*(\mathcal F)\bigr).
\]

% Source page: Exp1 component-local p. 44; combined-reader p. 54.
将后者记为 \({}^{t}\!\varphi_x(g)\). 则
\[
  {}^{t}\!\varphi_{gx}(h)\circ{}^{t}\!\varphi_x(g)
  =
  {}^{t}\!\bigl(\Lambda_{gx}(h)\circ\varphi_x(g)\bigr)^{-1}
  =
  \Lambda_x(hg)^{-1}
  =
  {}^{t}\!\varphi_x(hg).
\]
对每个 \(X\)-概形 \(f:Y\to X\), 有
\[
  V(\mathcal F)\times_fY\simeq V\bigl(f^*(\mathcal F)\bigr)
  \qquad\text{（参见 4.6.2）.}
\]
因此, 同构
\[
\begin{aligned}
  {}^{t}\!\varphi:\quad
  G\times_SV(\mathcal F)
  &=(G\times_SX)\times_{\mathrm{pr}_X}V(\mathcal F)
   =V\bigl(\mathrm{pr}_X^*(\mathcal F)\bigr)\\
  &\xrightarrow{\ \sim\ }
   V\bigl(\lambda^*(\mathcal F)\bigr)
   =(G\times_SX)\times_\lambda V(\mathcal F)
\end{aligned}
\]
赋予 \(V(\mathcal F)\) 一个 \(G\)-等变 \(\mathcal O_X\)-模结构.
将每个函子 \(V(-)\) 与表示它的向量纤维化等同, 再将
\({}^{t}\!\varphi\) 与到 \(V(\mathcal F)\) 的投影复合,
就得到 \(G\) 在 \(V(\mathcal F)\) 上的作用, 它与 \(X\) 上的作用相容,
且在各纤维上是线性的.

这样就得到了 [GIT, 第~1 章, \SGARefLink{sga3:i:section:3}{第~3 节}]
中的定义, 不同之处在于 Mumford 在那里考虑一个有限秩局部自由
\(\mathcal O_X\)-模 \(\mathcal E\), 并定义了 \(G\) 在
\(V(\mathcal E)=W(\mathcal E^\vee)\) 上的作用. 事实上, 将前图
\((\star\star)\) 中的 \(\theta\) 换成 \(\varphi\), 并反转各箭头的方向,
所得正是\emph{同上}定义~1.6 中的图; 上述同构
\({}^{t}\!\varphi\) 与\emph{同上}第~31 页的同构 \(\Phi\) 一致.

\medskip
\noindent\SGAInlineRefTarget{sga3:i:remark:6-5-4}\textbf{注 6.5.4.}
特别考虑 \(X=S\) 且 \(G\) 的作用平凡的情形.
此时, 一个 \(G\)-等变 \(\mathcal O_S\)-模 \(\mathcal F\)
就是一个 \(G\)-\(\mathcal O_S\)-模（参见 \SGARefLink{sga3:i:definition:4-7-1}{4.7.1}）.
若以 \(f:G\to S\) 表示此时同时等于 \(\mathrm{pr}_X\) 和
\(\lambda\) 的态射, 则同构
\[
  (\star)\qquad
  \theta:f^*(\mathcal F)\xrightarrow{\ \sim\ }f^*(\mathcal F)
\]
是下述集合的一个元素:
\[
\begin{aligned}
  \operatorname{Hom}_{\mathcal O_G}
    \bigl(f^*(\mathcal F),f^*(\mathcal F)\bigr)
  &=
  \operatorname{Hom}_{\mathcal O_G}
    \bigl(W(f^*(\mathcal F)),W(f^*(\mathcal F))\bigr)\\
  &=
  \operatorname{End}_{\mathcal O_S}\bigl(W(\mathcal F)\bigr)(G).
\end{aligned}
\]
它正是定义 \(G\) 在 \(W(\mathcal F)\) 上的作用的态射
\[
  \rho:G\longrightarrow
  \operatorname{End}_{\mathcal O_S}\bigl(W(\mathcal F)\bigr)
\]
. 此外, 由伴随关系, \(\theta\) 对应于 \(\mathcal O_S\)-模态射
\[
  f_*(\theta)\circ\tau:
  \mathcal F\longrightarrow f_*f^*(\mathcal F),
\]
其中 \(\tau:\mathcal F\to f_*f^*(\mathcal F)\) 是伴随关系的单位态射.
这一点将在第~VI\textsubscript{B} 篇的
\SGARefLink{cor:sga3-VIB-11.10}{11.10}~bis 中用到.

\subsection*{\SGAInlineRefTarget{sga3:i:subsection:6-6}6.6. 函子
\(\prod_{X/S}W(\mathcal F)\) 和 \(W(p_*(\mathcal F))\)}

设 \(S\) 是一个概形, \(G\) 是一个 \(S\)-群概形, 通过
\[
  \lambda:G\times_SX\to X,
  \qquad
  \mu:G\times_SY\to Y,
\]
作用在 \(S\)-概形 \(X\) 和 \(Y\) 上, 且 \(p:X\to Y\) 是一个
\(G\)-等变态射. 假设 \(p:X\to Y\) 和 \(\nu_Y:Y\to S\)
都是拟紧且拟分离的, 而 \(\pi:G\to S\) 是平坦的.

于是投影 \(\mathrm{pr}_Y:G\times_SY\to Y\) 是平坦的,
\(\mu\) 也是平坦的, 因为 \(\mu\) 是 \(\mathrm{pr}_Y\) 与自同构
\((g,y)\mapsto(g,gy)\) 的复合. 对变动的 \(S\)-概形
\(f:T\to S\), 以 \(p_T:X_T\to T\) 和 \(f_X:X_T\to X\)
表示由 \(p\) 和 \(f\) 作基变换所得的态射.

设 \(\mathcal F\) 是一个拟凝聚的 \(G\)-等变 \(\mathcal O_X\)-模,
并设
\[
  \theta:\mathrm{pr}_X^*(\mathcal F)
  \xrightarrow{\ \sim\ }\lambda^*(\mathcal F)
\]
是 \SGARefLink{sga3:i:definition:6-5-1}{6.5.1} \((\star)\) 中的同构.
由于 \(p\) 是拟紧且拟分离的, \(p_*(\mathcal F)\) 是拟凝聚的
（参见 EGA~I, 9.2.1）.

% Source page: Exp1 component-local p. 45; combined-reader p. 55.
\medskip
\noindent\SGAInlineRefTarget{sga3:i:lemma:6-6-1}\textbf{引理 6.6.1.}
\textit{拟凝聚 \(\mathcal O_Y\)-模 \(p_*(\mathcal F)\) 是 \(G\)-等变的.}

事实上, 有下列两个笛卡尔方块.

\begin{figure}[H]
  \centering
  \begin{tikzcd}[column sep=3.8em,row sep=3.2em]
  G\times_S X
    \arrow[r,"\mathrm{pr}_X"]
    \arrow[d,"q"']
  &
  X
    \arrow[d,"p"]
  &
  G\times_S X
    \arrow[l,"\lambda"']
    \arrow[d,"q"]
  \\
  G\times_S Y
    \arrow[r,"\mathrm{pr}_Y"]
  &
  Y
  &
  G\times_S Y
    \arrow[l,"\mu"']
  \qquad .
\end{tikzcd}

  \phantomsection\label{fig:sga3-exp1-p045-base-change-squares}
\end{figure}

由于 \(p\) 是拟紧且拟分离的, 而 \(\mathrm{pr}_Y\) 和 \(\mu\)
是平坦的, 由 EGA~III, 1.4.15 以及补充它的
EGA~IV\textsubscript{1}, 1.7.21, 得
\[
  q_*\mathrm{pr}_X^*(\mathcal F)
    =\mathrm{pr}_Y^*p_*(\mathcal F),
  \qquad
  q_*\lambda^*(\mathcal F)
    =\mu^*p_*(\mathcal F).
\]
所以 \(\theta\) 诱导一个同构
\[
  \theta_Y:\mathrm{pr}_Y^*p_*(\mathcal F)
  \xrightarrow{\ \sim\ }\mu^*p_*(\mathcal F).
\]
同样的论证表明, \(\theta_Y\) 满足余循环条件
\SGARefLink{sga3:i:definition:6-5-1}{6.5.1} \((\star\star)\),
因为所涉及的态射 \(G\times_SG\times_SY\to G\times_SY\) 都是平坦的.
引理得证.

特别地, 取 \(Y=S\), 其上 \(G\) 的作用为平凡作用.
由\SGARefLink{sga3:i:remark:6-5-4}{注~6.5.4}, 此时
\(p_*(\mathcal F)\) 是一个 \(G\)-\(\mathcal O_S\)-模.
此外, 若 \(G\) 在 \(S\) 上仿射, 且
\[
  \mathscr A(G)=\pi_*(\mathcal O_G),
\]
则由~\SGARefLink{sga3:i:proposition:4-7-2}{4.7.2},
\(p_*(\mathcal F)\) 是一个 \(\mathscr A(G)\)-余模.

另一方面, 根据\SGARefLink{sga3:i:proposition:6-4-1}{命题~6.4.1},
将每个 \(f:T\to S\) 映为
\[
  W\bigl(f_X^*(\mathcal F)\bigr)(X_T)
  =\Gamma\bigl(X_T,f_X^*(\mathcal F)\bigr)
  =\Gamma\bigl(T,p_{T*}f_X^*(\mathcal F)\bigr),
\]
的函子 \(\prod_{X/S}W(\mathcal F)\) 是一个
\(G\)-\(\mathcal O_S\)-模. 此外还有一个典范态射
\[
  \tau:W\bigl(p_*(\mathcal F)\bigr)
  \longrightarrow\prod_{X/S}W(\mathcal F),
\]
对每个 \(f:T\to S\), 它由典范态射
\[
  \Gamma\bigl(T,f^*p_*(\mathcal F)\bigr)
  \longrightarrow
  \Gamma\bigl(T,p_{T*}f_X^*(\mathcal F)\bigr).
\]
给出. 将其限制到由 \(S\) 上平坦的概形 \(T\) 组成的全子范畴,
就得到同构; 而 \(\tau\) 显然是 \(G\)-\(\mathcal O_S\)-模的态射.
于是我们证明了下列命题; 对于~(ii), 可比较 [GIT], 第~32 页.

\medskip
\noindent\SGAInlineRefTarget{sga3:i:proposition:6-6-2}\textbf{命题 6.6.2.}
\textit{设 \(S\) 是一个概形, \(G\) 是作用在 \(S\)-概形 \(X\)
上的一个 \(S\)-群概形, \(\mathcal F\) 是一个拟凝聚的 \(G\)-等变
\(\mathcal O_X\)-模. 假设 \(\pi:G\to S\) 是平坦的,
\(p:X\to S\) 是拟紧且拟分离的.}

\textit{(i) 则 \(p_*(\mathcal F)\) 是一个拟凝聚的
\(G\)-\(\mathcal O_S\)-模. 典范态射}
\[
  W\bigl(p_*(\mathcal F)\bigr)
  \longrightarrow\prod_{X/S}W(\mathcal F)
\]
\textit{是 \(G\)-\(\mathcal O_S\)-模的态射, 且这两个函子在平坦
\(S\)-概形范畴上相同.}

\textit{(ii) 此外, 若 \(G\) 在 \(S\) 上仿射, 且
\(\mathscr A(G)=\pi_*(\mathcal O_G)\), 则 \(p_*(\mathcal F)\)
具有一个 \(\mathscr A(G)\)-余模结构.}\footnote{编者注:
若 \(G\) 在 \(S\) 上平坦、拟紧且拟分离, 而 \(p_*(\mathcal F)\)
是一个平坦 \(\mathcal O_S\)-模, 则同样成立;
参见第~VI\textsubscript{B} 篇, \SGARefLink{concl:sga3-VIB-11.6.1-ii}{11.6.1(ii)}.}

% Source page: Exp1 component-local p. 46; combined-reader p. 56.
\subsection*{\SGAInlineRefTarget{sga3:i:subsection:6-7}6.7. 稳定子}

设 \(S\) 是一个概形, \(G\) 是一个作用在 \(S\)-概形 \(X\)
上的 \(S\)-群概形, \(\mathcal F\) 是一个拟凝聚的 \(G\)-等变
\(\mathcal O_X\)-模. 设 \(Y\) 是另一个带有 \(G\) 的作用的
\(S\)-概形, 该作用可以是平凡作用; 设 \(\phi:Y\to X\) 是一个
未必 \(G\)-等变的 \(S\)-态射, 并设
\(\operatorname{Stab}_G(\phi)\) 是 \(\phi\) 在 \(G\) 中的稳定子
（参见 \SGARefLink{sga3:i:definition:6-3-2}{6.3.2}）.

先指出（见第~VI\textsubscript{B} 篇,
\SGARefLink{sec:sga3-VIB-6}{第~6 节}）, 若 \(X\) 在 \(S\) 上分离,
且 \(Y\) 在 \(S\) 上本质自由, 则 \(\operatorname{Stab}_G(\phi)\)
可由 \(G\) 的闭子群概形 \(H\) 表示
（参见\emph{同上}, \SGARefLink{def:sga3-VIB-6.2.1}{定义~6.2.1}）.
事实上, 考虑
\[
  r:G\times_SY\longrightarrow X\times_SX,
  \qquad
  r(g,y)=\bigl(\phi(y),g\phi(g^{-1}y)\bigr),
\]
置 \(P=G\times_SY\), 并设 \(P'\) 是对角线 \(\Delta_{X/S}\)
在 \(r\) 下的逆像. 则（参见\emph{同上},
\SGARefLink{ex:sga3-VIB-6.2.4-a}{6.2.4(a)}）
\[
  \operatorname{Stab}_G(\phi)=\prod_{P/G}P'.
\]
由此可知, 若 \(X\) 在 \(S\) 上分离, 且 \(Y\) 在 \(S\) 上本质自由,
则 \(\operatorname{Stab}_G(\phi)\) 可由 \(G\) 的闭子群概形 \(H\) 表示.
后一条件在 \(S\) 是一个域的谱或 \(Y=S\) 时自动成立.
在这些假设下, \(\phi^*(\mathcal F)\) 是一个拟凝聚的
\(H\)-等变 \(\mathcal O_Y\)-模
（参见 \SGARefLink{sga3:i:remark:6-5-2}{6.5.2}）.

此外, 若 \(\pi:G\to S\) 和 \(p:Y\to S\) 在 \(S\) 上拟紧且拟分离,
且 \(\pi\) 平坦, 则由\SGARefLink{sga3:i:proposition:6-6-2}{命题~6.6.2},
\(p_*\phi^*(\mathcal F)\) 是一个 \(H\)-\(\mathcal O_S\)-模.
特别地:

\medskip
\noindent\SGAInlineRefTarget{sga3:i:corollary:6-7-1}\textbf{推论 6.7.1.}
\textit{设 \(S\) 是一个概形, \(G\) 是一个作用在 \(S\)-概形 \(X\)
上的 \(S\)-群概形, \(\mathcal F\) 是一个拟凝聚的 \(G\)-等变
\(\mathcal O_X\)-模. 假设 \(\pi:G\to S\) 平坦,
\(X\to S\) 拟紧且分离. 设 \(\tau:S\to X\) 是 \(X\) 在 \(S\) 上的一个截面.}

\textit{(i) 稳定子 \(H=\operatorname{Stab}_G(\tau)\) 是 \(G\) 的闭子群概形,
且 \(\tau^*(\mathcal F)\) 是一个拟凝聚的
\(H\)-\(\mathcal O_S\)-模.}

\textit{(ii) 若 \(H\) 还在 \(S\) 上仿射（例如 \(G\) 如此）, 且
\(\mathscr A(H)=\pi_*(\mathcal O_H)\), 则
\(\tau^*(\mathcal F)\) 是一个 \(\mathscr A(H)\)-余模.}

\subsection*{\SGAInlineRefTarget{sga3:i:subsection:6-8}6.8. \(G\) 上的 \(G\)-等变层}

最后给出两个结果 \SGARefLink{sga3:i:proposition:6-8-1}{6.8.1} 和
\SGARefLink{sga3:i:proposition:6-8-6}{6.8.6}, 它们将在第~II 篇与第~III 篇中用到,
特别是 III, \SGARefLink{sga3:iii:lemma:4-25}{4.25} 中.

\medskip
\noindent\SGAInlineRefTarget{sga3:i:proposition:6-8-1}\textbf{命题 6.8.1.}
\textit{设 \(S\) 是一个概形, \(\pi:G\to S\) 是一个 \(S\)-群概形,
\(\varepsilon:S\to G\) 是单位截面. 考虑 \(G\) 通过左平移在自身上的作用.
函子}
\[
  \mathcal E\longmapsto\pi^*(\mathcal E),
  \qquad
  \mathcal F\longmapsto\varepsilon^*(\mathcal F)
\]
\textit{在拟凝聚 \(\mathcal O_S\)-模范畴与拟凝聚的 \(G\)-等变
\(\mathcal O_G\)-模范畴之间诱导互为拟逆的等价.}

\textit{证明.}
以 \(\mu\) 表示 \(G\) 的乘法, 并设
\(\mathrm{pr}_2:G\times_SG\to G\) 是第二投影.
由于 \(\pi\circ\mu=\pi\circ\mathrm{pr}_2\), 对每个拟凝聚
\(\mathcal O_S\)-模 \(\mathcal E\), 都有典范同构
\[
  \mu^*\pi^*(\mathcal E)
  =\mathrm{pr}_2^*\pi^*(\mathcal E).
\]
此同构显然满足余循环条件
\SGARefLink{sga3:i:definition:6-5-1}{6.5.1} \((\star\star)\),
所以 \(\pi^*(\mathcal E)\) 是一个 \(G\)-等变 \(\mathcal O_G\)-模.
由于 \(\varepsilon^*\pi^*(\mathcal E)=\mathcal E\),
函子 \(\mathcal E\mapsto\pi^*(\mathcal E)\) 是

% Source page: Exp1 component-local p. 47; combined-reader p. 57.
全忠实的. 只须再证明, 对每个 \(G\)-等变 \(\mathcal O_G\)-模
\(\mathcal F\), 有
\[
  \mathcal F\simeq\pi^*\varepsilon^*(\mathcal F).
\]
由假设, 存在同构
\[
  \theta:\mathrm{pr}_2^*(\mathcal F)
  \xrightarrow{\ \sim\ }\mu^*(\mathcal F).
\]
将 \(\theta\) 沿态射
\[
  \tau:G\longrightarrow G\times_SG,
  \qquad
  \tau=(\operatorname{id}_G,\varepsilon\circ\pi),
\]
拉回, 就得到同构
\(\mathcal F\xrightarrow{\sim}\pi^*\varepsilon^*(\mathcal F)\).

\medskip
\noindent\SGAInlineRefTarget{sga3:i:remark:6-8-2}\textbf{注 6.8.2.}
\SGARefLink{sga3:i:definition:6-1:item:a}{(a)} 考虑由
\[
  (g_1,g_2)\cdot g=g_1gg_2^{-1}.
\]
定义的 \(G\times_SG\) 在 \(G\) 上的作用.
单位截面 \(\varepsilon:S\to G\) 的稳定子是
\(G\times_SG\) 的对角子群 \(H\).
因此, 若 \(\mathcal F\) 是一个拟凝聚的
\((G\times_SG)\)-等变 \(\mathcal O_G\)-模,
则由~\SGARefLink{sga3:i:remark:6-5-2}{6.5.2},
\(\mathcal E=\varepsilon^*(\mathcal F)\) 具有一个
\(H\)-\(\mathcal O_S\)-模结构.

(b) 可以证明, \(\mathcal F\mapsto\varepsilon^*(\mathcal F)\)
是拟凝聚的 \((G\times_SG)\)-等变 \(\mathcal O_G\)-模范畴与
拟凝聚的 \(H\)-\(\mathcal O_S\)-模范畴之间的一个等价.
这是更一般的下降结果的一个特例（参见第~IV 篇, 第~2 节, 以及 SGA~1, VIII）;
例如可见 [Th87], 1.2--1.3.

\medskip
\noindent\SGAInlineRefTarget{sga3:i:remark:6-8-3}\textbf{注 6.8.3.}
沿用 \SGARefLink{sga3:i:proposition:6-8-1}{6.8.1} 的记号.
对每个拟凝聚 \(\mathcal O_S\)-模 \(\mathcal E\),
以 \(\pi_*^G\pi^*(\mathcal E)\) 表示 \(\pi_*\pi^*(\mathcal E)\)
的 \(\mathcal O_S\)-子模, 它在 \(S\) 的开子集 \(V\) 上的截面
是所有满足如下条件的元素
\[
  \gamma\in\Gamma\bigl(\pi^{-1}(V),\pi^*(\mathcal E)\bigr)
\]
: 对每个 \(S'\to V\) 和 \(g\in G(S')\), 有
\(g\cdot\gamma_{S'}=\gamma_{S'}\). 则自然态射
\[
  \mathcal E\longrightarrow\pi_*^G\pi^*(\mathcal E)
\]
是同构. 若
\[
  \pi_*\pi^*(\mathcal E)
  =\mathcal E\otimes_{\mathcal O_S}\pi_*(\mathcal O_G),
\]
则这一点是直接的; 例如 \(G\to S\) 仿射时, 或 \(G\to S\)
拟紧且拟分离而 \(\mathcal E\) 平坦时就是如此.
一般情形也不难验证.

\medskip
\noindent\SGAInlineRefTarget{sga3:i:remark:6-8-4}\textbf{注 6.8.4.}
设 \(S\) 是一个概形, \(H\) 是作用在 \(S\)-概形 \(X\) 上的
一个 \(S\)-群概形, \(\mathcal F\) 是一个 \(\mathcal O_X\)-模.
假设 \(H\) 在 \(S\) 上平坦, 以 \(W_P(\mathcal F)\) 表示
\(W(\mathcal F)\) 在平坦 \(S\)-概形全子范畴上的限制.
由~\SGARefLink{sga3:i:definition:6-5-1}{6.5.1}, 赋予
\(\mathcal F\) 一个 \(H\)-等变模结构, 等价于给定
\(H\times_SX\) 上的一个满足余循环条件 \((\star\star)\) 的层同构
\[
  \theta:\mathrm{pr}_X^*(\mathcal F)\longrightarrow\lambda^*(\mathcal F)
\]
. 因此, 要赋予 \(\mathcal F\) 一个 \(H\)-等变模结构,
只须赋予 \(W_P(\mathcal F)\) 一个这样的结构.

\medskip
\SGAInlineRefTarget{sga3:i:recollection:6-8-5}\noindent\textbf{回顾 \SGARefLink{sga3:i:recollection:6-8-5}{6.8.5}.}
设 \(X\) 是一个 \(S\)-概形, \(Y\) 是 \(X\) 的一个闭
\(S\)-子概形. 以 \(\mathcal N_{Y/X}\) 表示浸入
\(i:Y\hookrightarrow X\) 的余法层（参见 EGA~IV\textsubscript{4},
16.1.2）. 若 \(S'\to S\) 平坦, 且
\(i':Y'\hookrightarrow X'\) 由 \(i\) 作基变换得到,
则由\emph{同上}, 16.2.2(iii), 得
\[
  \mathcal N_{Y/X}\otimes_{\mathcal O_Y}\mathcal O_{Y'}
  =\mathcal N_{Y'/X'}.
\]

\medskip
\noindent\SGAInlineRefTarget{sga3:i:proposition:6-8-6}\textbf{命题 6.8.6.}
\textit{设 \(S\) 是一个概形, \(X\) 是一个 \(S\)-群概形,
\(Y\) 是 \(X\) 的一个 \(S\)-子群概形. 假设 \(Y\) 在 \(S\)
上平坦. 则余法层 \(\mathcal N_{Y/X}\) 是一个
\((Y\times_SY)\)-等变 \(\mathcal O_Y\)-模.}

事实上, \(H=Y\times_SY\) 在 \(S\) 上平坦.
因此, 根据\SGARefLink{sga3:i:remark:6-8-4}{注~6.8.4},
只须赋予 \(W_P(\mathcal N_{Y/X})\) 一个 \(H\)-等变模结构.
设 \(S'\) 是一个平坦 \(S\)-概形, \(Y'\hookrightarrow X'\)
是作基变换得到的浸入, 并置
\[
  \mathcal N'
  =\mathcal N_{Y/X}\otimes_{\mathcal O_Y}\mathcal O_{Y'}.
\]
根据~\SGARefLink{sga3:i:recollection:6-8-5}{6.8.5},
\(\mathcal N'=\mathcal N_{Y'/X'}\).
每个 \(h\in H(S')\) 都诱导 \(Y'\) 的一个自同构,
从而对每个 \(y\in Y'(S')\), 诱导同构
\[
  \Gamma\bigl(S',y^*(\mathcal N')\bigr)
  \xrightarrow{\ \sim\ }
  \Gamma\bigl(S',y^*h^*(\mathcal N')\bigr)
\]

% Source page: Exp1 component-local p. 48; combined-reader p. 58.
这些同构赋予 \(W_P(\mathcal N_{Y/X})\) 一个 \(H\)-等变模结构
（参见 \SGARefLink{sga3:i:definition:6-1}{6.1}）.

\section*{参考文献}
\addcontentsline{toc}{section}{参考文献}
\footnote{编者注: 这些是本篇所引用的补充参考文献.}

\begin{description}
\item[{[DG70]}] M. Demazure 和 P. Gabriel,
  \emph{Groupes Alg\'ebriques}, Masson \& North-Holland, 1970.
\item[{[Gr57]}] A. Grothendieck,
  "Sur quelques points d'alg\`ebre homologique,"
  \emph{T\^ohoku Math. J.} \textbf{9} (1957), 119--221.
\item[{[Ja03]}] J. C. Jantzen,
  \emph{Representations of Algebraic Groups}, Academic Press, 1987;
  第 2 版, American Mathematical Society, 2003.
\item[{[GIT]}] D. Mumford,
  \emph{Geometric Invariant Theory}, Springer-Verlag, 1965;
  与 J. Fogarty 合著的第 2 版, 1982; 与 J. Fogarty
  和 F. Kirwan 合著的第 3 版, 1994.
\item[{[Ni02]}] N. Nitsure,
  "Representability of \(\mathrm{GL}_{\mathcal E}\),"
  \emph{Proc. Indian Acad. Sci.} \textbf{112} (2002), 第~4 期, 539--542.
\item[{[Ni04]}] N. Nitsure,
  "Representability of Hom implies flatness,"
  \emph{Proc. Indian Acad. Sci.} \textbf{114} (2004), 第~1 期, 7--14.
\item[{[Se68]}] J.-P. Serre,
  "Groupes de Grothendieck des sch\'emas en groupes r\'eductifs
  d\'eploy\'es,"
  \emph{Publ. Math. I.H.\'E.S.} \textbf{34} (1968), 37--52.
\item[{[Th87]}] R. W. Thomason,
  "Equivariant resolution, linearization, and Hilbert's fourteenth
  problem over arbitrary base schemes,"
  \emph{Adv. Math.} \textbf{65} (1987), 16--34.
\end{description}

% End of Expose I and of the assigned component.

% Translation checkpoint: source lines 881; translated lines 849 (excluding this statistics comment).
% Math segments: 651 (source and translation); all math tokens preserved except two translated \text prose arguments.
% Unresolved source-verification questions: 6.5.3 transpose equality; 6.6 general-base p_T target. Formulas retained verbatim.
